\documentclass[12pt,a4paper]{article}

% ---- Encoding & Font ----
\usepackage{iftex}
\ifPDFTeX
  \usepackage[utf8]{inputenc}
  \usepackage[T1]{fontenc}
\else
  \usepackage{fontspec}  % XeTeX/LuaTeX (tectonic): Unicode nativo (·, —, ¿, ª, §)
\fi
\usepackage[spanish]{babel}
\usepackage{csquotes}

% ---- Page layout ----
\usepackage[top=2.5cm, bottom=2.5cm, left=2.5cm, right=2.5cm]{geometry}
\usepackage{setspace}
\onehalfspacing

% ---- Math ----
\usepackage{amsmath, amssymb, amsthm}
\usepackage{mathtools}
\usepackage{physics}

% ---- Graphics ----
\usepackage{graphicx}
\usepackage{tikz}
\usepackage{float}
\usepackage{caption}

% ---- Tables ----
\usepackage{array}
\usepackage{booktabs}
\usepackage{tabularx}
\usepackage{colortbl}

% ---- Colours ----
\usepackage{xcolor}
\definecolor{notebg}{HTML}{FFF3CD}
\definecolor{noteborder}{HTML}{FFC107}
\definecolor{boxred}{HTML}{F8D7DA}
\definecolor{boxredborder}{HTML}{F5C6CB}

% ---- Boxes ----
\usepackage{mdframed}
\usepackage{tcolorbox}
\tcbuselibrary{most}

\newtcolorbox{keybox}[1][]{
  colback=blue!5,
  colframe=blue!70!black,
  arc=4pt,
  boxrule=1pt,
  fonttitle=\bfseries,
  title={#1}
}

\newtcolorbox{warnbox}[1][]{
  colback=boxred,
  colframe=boxredborder,
  coltitle=black!75,
  arc=4pt,
  boxrule=1pt,
  fonttitle=\bfseries,
  title={#1}
}

\newtcolorbox{notebox}[1][]{
  colback=notebg,
  colframe=noteborder,
  coltitle=black!75,
  arc=4pt,
  boxrule=1pt,
  fonttitle=\bfseries,
  title={#1}
}

% ---- Hyperlinks & TOC ----
\usepackage[hidelinks]{hyperref}
\usepackage{bookmark}

% ---- Enumerate ----
\usepackage{enumitem}

% ---- Miscellaneous ----
\usepackage{icomma}
\usepackage{siunitx}
\usepackage{textcomp}
\usepackage[bottom]{footmisc}

% ---- Pseudocódigo ----
\usepackage[spanish,onelanguage,ruled,vlined]{algorithm2e}
\SetKwInput{KwEntrada}{Entrada}
\SetKwInput{KwSalida}{Salida}

% ---- Diagramas (gráficas) ----
\usepackage{pgfplots}
\pgfplotsset{compat=1.18}
\usetikzlibrary{babel}  % babel español activa `>`; esto evita romper las flechas `->`
% courses/MetNum2023/Clases/assets/tikzstyles.tex
% ---------------------------------------------------------------------------
% Estilos compartidos para los diagramas del curso Métodos Numéricos
% (MetNum2023). Fuente única del "look" visual (colores, grosores, escalas)
% para que todas las figuras (matrices triangulares, curvas de convergencia
% / error, splines a trozos) sean consistentes entre clases.
%
% Uso: la clase debe cargar antes `tikz` y `pgfplots` (bloque §6.3 del
%      CLAUDE.md del curso), y luego
%      % courses/MetNum2023/Clases/assets/tikzstyles.tex
% ---------------------------------------------------------------------------
% Estilos compartidos para los diagramas del curso Métodos Numéricos
% (MetNum2023). Fuente única del "look" visual (colores, grosores, escalas)
% para que todas las figuras (matrices triangulares, curvas de convergencia
% / error, splines a trozos) sean consistentes entre clases.
%
% Uso: la clase debe cargar antes `tikz` y `pgfplots` (bloque §6.3 del
%      CLAUDE.md del curso), y luego
%      % courses/MetNum2023/Clases/assets/tikzstyles.tex
% ---------------------------------------------------------------------------
% Estilos compartidos para los diagramas del curso Métodos Numéricos
% (MetNum2023). Fuente única del "look" visual (colores, grosores, escalas)
% para que todas las figuras (matrices triangulares, curvas de convergencia
% / error, splines a trozos) sean consistentes entre clases.
%
% Uso: la clase debe cargar antes `tikz` y `pgfplots` (bloque §6.3 del
%      CLAUDE.md del curso), y luego
%      \input{../assets/tikzstyles.tex}
% (compilar desde el directorio de la clase para que la ruta relativa resuelva).
%
% Paleta propia de este curso (no se reusa la de Fisica3-2015/CDIV2017/
% ElecMag2024): mismos tonos base para alinear con keybox/notebox/warnbox,
% pero archivo independiente, sin dependencia cruzada entre cursos.
% ---------------------------------------------------------------------------

% ---- Paleta (alineada con las cajas del preámbulo: keybox/notebox/warnbox) ----
\definecolor{figblue}{HTML}{1F4E79}   % azul principal (L, curvas principales, keybox)
\definecolor{figred}{HTML}{B03A2E}    % rojo de énfasis (U, warnbox, error)
\definecolor{figamber}{HTML}{B8860B}  % ámbar (notebox, pivotes)
\definecolor{figgray}{HTML}{555555}   % auxiliares, ejes, ceros de matriz

% ---- TikZ: librerías y estilos reutilizables ----
% Nota: las puntas se declaran como `-{Latex[...]}` y no como `->`. La forma
% con llaves NO contiene un `>` literal, así que es inmune al `>` activo que
% introduce babel español (de todos modos se carga \usetikzlibrary{babel} en
% el bloque §6.3 para cubrir cualquier `->` suelto).
\usetikzlibrary{arrows.meta,calc,matrix,patterns}

\tikzset{
  figline/.style={figblue, line width=0.9pt},
  figaux/.style={figgray, dashed, line width=0.6pt},
  figaxis/.style={figgray, line width=0.7pt, -{Latex[length=2.2mm,width=1.5mm]}},
  % estructura de matrices (L triangular inferior, U triangular superior)
  matcelda/.style={minimum width=7mm, minimum height=7mm, anchor=center, font=\small},
  matnz/.style={matcelda, fill=figblue!12},
  matnzU/.style={matcelda, fill=figred!12},
  matcero/.style={matcelda, text=figgray!70},
  matpivote/.style={matcelda, fill=figamber!25, draw=figamber, line width=0.8pt},
  % nodos de interpolación / splines
  fignodo/.style={circle, inner sep=0pt, minimum size=4.5pt, draw=figblue, line width=0.9pt, fill=white},
  fignodoactivo/.style={fignodo, fill=figamber},
  % etiquetas
  figlbl/.style={font=\small},
  figlblaux/.style={font=\footnotesize, figgray},
}

% ---- pgfplots: estilos base ----
\pgfplotsset{
  % convergencia / error vs. n (grado, iteración) — suele ir en escala log
  errorfig/.style={
    width=10cm, height=6.5cm,
    xlabel style={font=\small},
    ylabel style={font=\small},
    tick label style={font=\footnotesize},
    every axis plot/.append style={line width=1.1pt, mark=*, mark size=1.5pt},
    grid=major, grid style={figgray!15},
    legend cell align=left,
    legend style={font=\footnotesize, draw=figgray!40, at={(0.5,-0.25)}, anchor=north},
    clip=false,
  },
  % interpolación / splines: un color por tramo, nodos marcados aparte
  interpfig/.style={
    width=10cm, height=6cm,
    axis lines=middle,
    xlabel style={font=\small, at={(axis description cs:1,0.5)}, anchor=west},
    ylabel style={font=\small, at={(axis description cs:0.5,1)}, anchor=south},
    tick label style={font=\footnotesize},
    every axis plot/.append style={line width=1.1pt},
    grid=none,
    legend cell align=left,
    legend style={font=\footnotesize, draw=figgray!40},
    clip=false,
  },
}

% (compilar desde el directorio de la clase para que la ruta relativa resuelva).
%
% Paleta propia de este curso (no se reusa la de Fisica3-2015/CDIV2017/
% ElecMag2024): mismos tonos base para alinear con keybox/notebox/warnbox,
% pero archivo independiente, sin dependencia cruzada entre cursos.
% ---------------------------------------------------------------------------

% ---- Paleta (alineada con las cajas del preámbulo: keybox/notebox/warnbox) ----
\definecolor{figblue}{HTML}{1F4E79}   % azul principal (L, curvas principales, keybox)
\definecolor{figred}{HTML}{B03A2E}    % rojo de énfasis (U, warnbox, error)
\definecolor{figamber}{HTML}{B8860B}  % ámbar (notebox, pivotes)
\definecolor{figgray}{HTML}{555555}   % auxiliares, ejes, ceros de matriz

% ---- TikZ: librerías y estilos reutilizables ----
% Nota: las puntas se declaran como `-{Latex[...]}` y no como `->`. La forma
% con llaves NO contiene un `>` literal, así que es inmune al `>` activo que
% introduce babel español (de todos modos se carga \usetikzlibrary{babel} en
% el bloque §6.3 para cubrir cualquier `->` suelto).
\usetikzlibrary{arrows.meta,calc,matrix,patterns}

\tikzset{
  figline/.style={figblue, line width=0.9pt},
  figaux/.style={figgray, dashed, line width=0.6pt},
  figaxis/.style={figgray, line width=0.7pt, -{Latex[length=2.2mm,width=1.5mm]}},
  % estructura de matrices (L triangular inferior, U triangular superior)
  matcelda/.style={minimum width=7mm, minimum height=7mm, anchor=center, font=\small},
  matnz/.style={matcelda, fill=figblue!12},
  matnzU/.style={matcelda, fill=figred!12},
  matcero/.style={matcelda, text=figgray!70},
  matpivote/.style={matcelda, fill=figamber!25, draw=figamber, line width=0.8pt},
  % nodos de interpolación / splines
  fignodo/.style={circle, inner sep=0pt, minimum size=4.5pt, draw=figblue, line width=0.9pt, fill=white},
  fignodoactivo/.style={fignodo, fill=figamber},
  % etiquetas
  figlbl/.style={font=\small},
  figlblaux/.style={font=\footnotesize, figgray},
}

% ---- pgfplots: estilos base ----
\pgfplotsset{
  % convergencia / error vs. n (grado, iteración) — suele ir en escala log
  errorfig/.style={
    width=10cm, height=6.5cm,
    xlabel style={font=\small},
    ylabel style={font=\small},
    tick label style={font=\footnotesize},
    every axis plot/.append style={line width=1.1pt, mark=*, mark size=1.5pt},
    grid=major, grid style={figgray!15},
    legend cell align=left,
    legend style={font=\footnotesize, draw=figgray!40, at={(0.5,-0.25)}, anchor=north},
    clip=false,
  },
  % interpolación / splines: un color por tramo, nodos marcados aparte
  interpfig/.style={
    width=10cm, height=6cm,
    axis lines=middle,
    xlabel style={font=\small, at={(axis description cs:1,0.5)}, anchor=west},
    ylabel style={font=\small, at={(axis description cs:0.5,1)}, anchor=south},
    tick label style={font=\footnotesize},
    every axis plot/.append style={line width=1.1pt},
    grid=none,
    legend cell align=left,
    legend style={font=\footnotesize, draw=figgray!40},
    clip=false,
  },
}

% (compilar desde el directorio de la clase para que la ruta relativa resuelva).
%
% Paleta propia de este curso (no se reusa la de Fisica3-2015/CDIV2017/
% ElecMag2024): mismos tonos base para alinear con keybox/notebox/warnbox,
% pero archivo independiente, sin dependencia cruzada entre cursos.
% ---------------------------------------------------------------------------

% ---- Paleta (alineada con las cajas del preámbulo: keybox/notebox/warnbox) ----
\definecolor{figblue}{HTML}{1F4E79}   % azul principal (L, curvas principales, keybox)
\definecolor{figred}{HTML}{B03A2E}    % rojo de énfasis (U, warnbox, error)
\definecolor{figamber}{HTML}{B8860B}  % ámbar (notebox, pivotes)
\definecolor{figgray}{HTML}{555555}   % auxiliares, ejes, ceros de matriz

% ---- TikZ: librerías y estilos reutilizables ----
% Nota: las puntas se declaran como `-{Latex[...]}` y no como `->`. La forma
% con llaves NO contiene un `>` literal, así que es inmune al `>` activo que
% introduce babel español (de todos modos se carga \usetikzlibrary{babel} en
% el bloque §6.3 para cubrir cualquier `->` suelto).
\usetikzlibrary{arrows.meta,calc,matrix,patterns}

\tikzset{
  figline/.style={figblue, line width=0.9pt},
  figaux/.style={figgray, dashed, line width=0.6pt},
  figaxis/.style={figgray, line width=0.7pt, -{Latex[length=2.2mm,width=1.5mm]}},
  % estructura de matrices (L triangular inferior, U triangular superior)
  matcelda/.style={minimum width=7mm, minimum height=7mm, anchor=center, font=\small},
  matnz/.style={matcelda, fill=figblue!12},
  matnzU/.style={matcelda, fill=figred!12},
  matcero/.style={matcelda, text=figgray!70},
  matpivote/.style={matcelda, fill=figamber!25, draw=figamber, line width=0.8pt},
  % nodos de interpolación / splines
  fignodo/.style={circle, inner sep=0pt, minimum size=4.5pt, draw=figblue, line width=0.9pt, fill=white},
  fignodoactivo/.style={fignodo, fill=figamber},
  % etiquetas
  figlbl/.style={font=\small},
  figlblaux/.style={font=\footnotesize, figgray},
}

% ---- pgfplots: estilos base ----
\pgfplotsset{
  % convergencia / error vs. n (grado, iteración) — suele ir en escala log
  errorfig/.style={
    width=10cm, height=6.5cm,
    xlabel style={font=\small},
    ylabel style={font=\small},
    tick label style={font=\footnotesize},
    every axis plot/.append style={line width=1.1pt, mark=*, mark size=1.5pt},
    grid=major, grid style={figgray!15},
    legend cell align=left,
    legend style={font=\footnotesize, draw=figgray!40, at={(0.5,-0.25)}, anchor=north},
    clip=false,
  },
  % interpolación / splines: un color por tramo, nodos marcados aparte
  interpfig/.style={
    width=10cm, height=6cm,
    axis lines=middle,
    xlabel style={font=\small, at={(axis description cs:1,0.5)}, anchor=west},
    ylabel style={font=\small, at={(axis description cs:0.5,1)}, anchor=south},
    tick label style={font=\footnotesize},
    every axis plot/.append style={line width=1.1pt},
    grid=none,
    legend cell align=left,
    legend style={font=\footnotesize, draw=figgray!40},
    clip=false,
  },
}


% ---- Title ----
\title{\textbf{Clase 6: Descomposición LU} \\[0.3em]
        \large{Pivoteo parcial, matrices de multiplicadores y el teorema $PA=LU$}}
\author{Transcripto y expandido --- Curso Métodos Numéricos (OpenFING)}
\date{}

\begin{document}

\maketitle
\thispagestyle{empty}
\newpage

\tableofcontents
\newpage

\section{Motivación: por qué falla la eliminación gaussiana sin pivoteo}

Retomamos un ejemplo de la clase anterior: un sistema $Ax=b$ que \textbf{no
tiene ninguna dificultad especial}, resuelto con eliminación gaussiana
\textbf{sin pivoteo} en una máquina hipotética de 5 cifras significativas. El
resultado numérico queda muy lejos de la solución real.

Diagnóstico paso a paso:

\begin{itemize}
  \item En el segundo paso de la eliminación aparece un multiplicador
  $L_{32} \approx 2500$ — mucho más grande que el orden de magnitud de los
  coeficientes del sistema (del orden de $1$).
  \item Al calcular $b_3 \leftarrow b_3 + L_{32}\cdot(\text{fila }2)$, ese
  multiplicador gigante amplifica cualquier error relativo pequeño (de
  redondeo, inevitable en 5 cifras) hasta convertirlo en un error
  \textbf{absoluto grande} frente a la magnitud real de los coeficientes.
  \item El error se arrastra a la sustitución hacia atrás: al despejar $x_2$
  (multiplicado por un coeficiente chico), un error absoluto que ya era
  grande se magnifica todavía más.
\end{itemize}

\begin{notebox}[Idea clave]
No es que \enquote{dividir} o el pivote en sí sean malos — la división por un
pivote no introduce error relativo grande. El problema es puramente de
\textbf{magnitud de los multiplicadores}: si $L_{jk}$ es grande, cualquier
error de redondeo en los datos se amplifica al multiplicarlo.
\end{notebox}

\[
\boxed{\text{Multiplicadores grandes} \implies \text{amplificación de error de redondeo}}
\]

\section{Pivoteo parcial}

\subsection{Estrategia}

En el paso $k$-ésimo, en vez de usar $A_{kk}$ como pivote sin más, se mira
toda la columna $k$ de la parte no reducida (filas $k,\dots,n$) y se elige
como pivote el elemento de mayor valor absoluto:

\[
p = \operatorname*{argmax}_{j \geq k} |A_{jk}|
\]

(si hay empate, cualquier índice que alcance el máximo sirve). Se intercambia
la fila $k$ con la fila $p$. Esto acota todos los multiplicadores:

\[
\boxed{|L_{jk}| \leq 1 \quad \text{para todo } j > k}
\]

lo que evita exactamente el problema de la Sección 1.

\begin{notebox}[Alcance del curso]
Solo se trabaja con pivoteo \textbf{de filas} (pivoteo parcial). Existe
también pivoteo de columnas o pivoteo total, pero no se desarrolla: alcanza
y sobra para los efectos del curso.
\end{notebox}

\subsection{Costo adicional}

Elegir el pivote en el paso $k$ cuesta del orden de $n-k-1$ comparaciones.
Sumando sobre todos los pasos:

\[
\sum_{k=1}^{n-1} (n-k-1) = O(n^2)
\]

frente al costo de la eliminación en sí, $O(n^3)$ (más precisamente
$\tfrac{2}{3}n^3$).

\begin{keybox}[Costo de pivotear]
El costo de eliminación gaussiana con pivoteo parcial es comparable al de la
eliminación sin pivoteo: no se paga mucho más por la seguridad adicional. El
intercambio de filas en sí también es $O(n^2)$ en total (lineal en $n$ por
cada uno de los $n-1$ pasos) — tampoco cambia el orden dominante.
\end{keybox}

\section{Codificar la eliminación como producto de matrices}

La eliminación gaussiana con pivoteo parcial —vista antes como una receta de
pasos— se reinterpreta como una sucesión de multiplicaciones de $A$ por
matrices especiales. Esto es lo que permite después probar el teorema de la
descomposición LU.

\subsection{Matrices de permutación}

Una \textbf{matriz de permutación} $P$ se obtiene intercambiando filas (o
columnas) de la identidad. Multiplicar por una matriz de permutación permuta
filas o columnas:

\begin{itemize}
  \item $P \cdot A$ permuta \textbf{filas} de $A$.
  \item $A \cdot P$ permuta \textbf{columnas} de $A$.
\end{itemize}

El producto de matrices de permutación es una matriz de permutación.

\begin{notebox}[Detalle de implementación]
Una matriz de permutación $n\times n$ tiene $n^2$ entradas pero solo $n$ son
información (dónde están los unos). En Octave conviene guardar un
\textbf{vector de permutación}: por ejemplo, intercambiar las filas 1 y 3 de
una identidad $4\times 4$ se guarda como el vector $(3,2,1,4)$, y entonces
$P\cdot A$ se escribe \texttt{A(p,:)} y $A \cdot P$ se escribe \texttt{A(:,p)}.
\end{notebox}

\subsection{Matrices de multiplicadores}

El paso $k$ de la eliminación (restar a cada fila $j>k$ un múltiplo $L_{jk}$
de la fila $k$) se escribe como multiplicar $A$ por una matriz $M_k$:

\begin{itemize}
  \item $M_k$ es la identidad, salvo que en la \textbf{columna $k$}, por
  debajo de la diagonal, tiene las entradas $-L_{k+1,k},\dots,-L_{n,k}$.
  \item Es decir: $M_k$ es triangular inferior, con unos en la diagonal, y
  todos sus elementos no nulos fuera de la diagonal están en \textbf{una
  sola columna}.
\end{itemize}

\begin{figure}[H]
\begin{figure}[H]\centering
\begin{tikzpicture}
\matrix (a) [matrix of math nodes,left delimiter=(,right delimiter=),nodes={matcelda},row sep=2pt,column sep=4pt] {
1&0&0&0\\
-\ell_{21}&1&0&0\\
-\ell_{31}&0&1&0\\
-\ell_{41}&0&0&1\\
};
\draw[figamber,thick,rounded corners] ($(a-2-1.north west)+(-0.08,0.08)$) rectangle ($(a-4-1.south east)+(0.08,-0.08)$);
\node[anchor=west,font=\small] (t) at (4,0) {$F_j\gets F_j-\ell_{j1}F_1$};
\draw[figaux] (a-3-1.west)--++(-0.7,0) node[left,font=\small] {columna 1};
\node[font=\small] at (0,-2.4) {$M_1$: diagonal de unos};
\end{tikzpicture}
\caption{La primera matriz de eliminación almacena los negativos de los multiplicadores en una sola columna. Multiplicarla por $A$ a la izquierda efectúa simultáneamente las combinaciones de filas; invertirla cambia esos signos.}
\end{figure}

\caption{Estructura de una matriz de multiplicadores $M_k$ ($n=5$, $k=3$):
identidad salvo la columna $k$ por debajo de la diagonal, donde aparecen los
multiplicadores $-L_{j,k}$.}
\end{figure}

Dos propiedades, verificables \enquote{mirando} las matrices (sin hacer
cuentas), que son la clave técnica de la demostración:

\begin{enumerate}
  \item El producto de dos matrices de multiplicadores de pasos distintos se
  obtiene \textbf{pegando} las columnas no nulas de cada una.
  \item La inversa de una matriz de multiplicadores es la misma matriz con
  el signo de los multiplicadores cambiado (invertir, en general, cuesta
  resolver $n$ sistemas; acá es trivial).
\end{enumerate}

\subsection{La relación de conmutación \texorpdfstring{$PM = \tilde{M}P$}{PM = M-tilde P}}

Cuando una matriz de permutación \textbf{de un paso posterior} multiplica a
una matriz de multiplicadores de un paso anterior, se puede pasar la
permutación al otro lado, a costa de permutar también las entradas de la
matriz de multiplicadores:

\[
P \cdot M = \tilde{M} \cdot P
\]

donde $\tilde{M}$ tiene la misma estructura que $M$, solo que con los
multiplicadores reordenados según $P$. Esta relación permite, más adelante,
separar todas las $P$ de un lado y todas las $M$ del otro en la cadena de
productos.

\section{El teorema de descomposición LU}

\subsection{Demostración (constructiva)}

La eliminación gaussiana con pivoteo parcial, vista como producto de
matrices, se escribe (para $A$ de $n\times n$):

\[
M_{n-1} P_{n-1} \cdots M_2 P_2\, M_1 P_1\, A = U
\]

donde cada $P_k$ es la matriz de permutación del paso $k$ (la identidad si
no hubo que pivotear) y cada $M_k$ es la matriz de multiplicadores de ese
paso. $U$ queda triangular superior (aunque $A$ sea singular, se llega
igual a una triangular superior, eventualmente con algún cero en la
diagonal).

Usando la asociatividad del producto de matrices y la relación de
conmutación de la Sección 3.3 repetidamente, se agrupan todas las $P_k$ a la
izquierda (formando una única $P$) y todas las $M_k$ transformadas del otro
lado. Pasando las matrices de multiplicadores al otro lado de la igualdad
(invirtiéndolas, lo cual es solo cambiar signos) y multiplicándolas entre sí
(pegando columnas), se llega a:

\[
P\,A = L\,U
\]

donde $P$ es de permutación, $U$ es la triangular superior de la
eliminación, y $L$ es triangular inferior con unos en la diagonal, cuyas
entradas fuera de la diagonal son literalmente los multiplicadores $L_{jk}$
de la eliminación (con el efecto de las permutaciones incorporado).

\begin{keybox}[Teorema de descomposición LU]
Para toda matriz cuadrada $A$, existen $P$ (permutación), $L$ (triangular
inferior con unos en la diagonal) y $U$ (triangular superior) tales que
\[
PA = LU.
\]
\end{keybox}

\subsection{Comentarios sobre el enunciado}

\begin{itemize}
  \item \textbf{La $P$ es necesaria en general}: no toda matriz cuadrada se
  puede escribir como $A=LU$ sin permutar filas.
  \item \textbf{$\det A$ puede ser cualquiera}: el teorema no exige
  $\det A \neq 0$. Si $A$ es singular, $U$ sale con algún cero en la
  diagonal (y eso revela que $A$ es singular). El curso trabaja
  mayormente con $A$ no singular porque interesa resolver sistemas, pero el
  teorema en sí no lo necesita.
  \item \textbf{Costo}: como la demostración es constructiva y ligada al
  algoritmo, calcular $P,L,U$ cuesta como mucho lo mismo que la eliminación
  con pivoteo parcial: $O(n^3)$ (del orden de $\tfrac{2}{3}n^3$).
\end{itemize}

\subsection{No unicidad}

La factorización $PA=LU$ \textbf{no es única}: depende de qué pivoteo se
hizo (libertad al resolver empates en el máximo), y aun fijando \enquote{unos
en la diagonal de $L$} (ya una elección), un empate entre candidatos a
pivote deja libertad de pivotear o no, lo que cambia $P,L,U$.

\begin{notebox}
No confundir con la factorización de Cholesky (Sección 7), que sí es única
bajo sus hipótesis — son construcciones distintas.
\end{notebox}

\section{Resolver sistemas usando la LU}

Dada $PA=LU$ con $A$ no singular, resolver $Ax=b$ se reduce a dos sistemas
triangulares: multiplicar por $P$ ($PAx=Pb \Rightarrow LUx=Pb$), definir
$y:=Ux$, y resolver $Ly=Pb$ (sustitución hacia adelante) seguido de $Ux=y$
(sustitución hacia atrás).

\begin{algorithm}[H]
\SetAlgoLined
\KwEntrada{$L,U$ triangulares, vector de permutación $p$, $b \in \mathbb{R}^n$}
\KwSalida{$x$ tal que $Ax=b$}
$b' \gets b(p)$ \tcp*{reordenar $b$ según la permutación}
\tcp{sustitución hacia adelante: $Ly=b'$}
\For{$i \gets 1$ \KwTo $n$}{
  $y_i \gets \left(b'_i - \sum_{j=1}^{i-1} L_{ij}\, y_j\right)$\;
}
\tcp{sustitución hacia atrás: $Ux=y$}
\For{$i \gets n, n-1, \dots, 1$}{
  $x_i \gets \left(y_i - \sum_{j=i+1}^{n} U_{ij}\, x_j\right) / U_{ii}$\;
}
\caption{Resolver $Ax=b$ a partir de una factorización $PA=LU$ ya calculada}
\end{algorithm}

Cada sustitución cuesta $O(n^2)$.

\[
\boxed{\text{Con } P,L,U \text{ ya calculados, resolver } Ax=b \text{ cuesta } O(n^2), \text{ no } O(n^3)}
\]

\section{Cuándo conviene calcular la LU}

Calcular $P,L,U$ cuesta $O(n^3)$, el mismo orden que resolver el sistema una
vez por eliminación directa. La ventaja aparece cuando hay que resolver
$Ax=b$ para \textbf{muchos $b$ distintos} con la misma $A$.

\textbf{Ejemplo (ingeniería civil):} un puente modelado por $Ax=b$, donde
$A$ depende de la estructura (fija) y $b$ representa distintas
configuraciones de carga. Si hay $m$ configuraciones a evaluar:

\begin{table}[H]
\centering
\begin{tabular}{@{}ll@{}}
\toprule
\textbf{Estrategia} & \textbf{Costo total} \\
\midrule
Eliminación gaussiana completa para cada $b$ & $O(n^3 \cdot m)$ \\
Calcular $P,L,U$ una vez + sustitución para cada $b$ & $O(n^3 + n^2 \cdot m)$ \\
\bottomrule
\end{tabular}
\end{table}

Si $m$ es comparable a $n$: $O(n^4)$ contra $O(n^3)$ — se gana un orden de
magnitud completo.

\begin{keybox}[Regla práctica]
Si la matriz se reutiliza muchas veces con distintos lados derechos,
conviene factorizar una sola vez y reutilizar $L$ y $U$.
\end{keybox}

\section{Nota al margen: la factorización de Cholesky}

Si $A$ es simétrica ($A=A^T$), por el teorema espectral es diagonalizable en
$\mathbb{R}$ con $n$ valores propios reales. Si además todos son $\geq 0$,
$A$ se dice semidefinida positiva (definida positiva si son estrictamente
$>0$). Para esas matrices existe una factorización distinta de la LU:

\[
A = R^T R
\]

con $R$ triangular (\textbf{Cholesky}) — no tiene unos en la diagonal, no
tiene que ver con eliminación gaussiana, y es un poco más barata de calcular
que la LU.

\begin{notebox}
No confundir con la descomposición LU de la Sección 4: son construcciones
distintas, con hipótesis distintas sobre $A$. No se desarrolla en esta
clase (queda para el práctico).
\end{notebox}

\section{El operador backslash en MATLAB/Octave}

El operador \textbf{backslash} (\texttt{x = A \textbackslash\ b}) decide
internamente qué método usar según la estructura de $A$. La clase muestra
una versión \enquote{de juguete} tomada del libro de \textbf{Cleve Moler},
uno de los fundadores de MATLAB —autor de uno de los libros de referencia
del curso, con muchos ejemplos de código (el ejemplo de la máquina de 5
dígitos de la Sección 1 sale de ese mismo libro).

\begin{notebox}[Nota de transcripción]
El audio transcribe el nombre del autor como \enquote{Cliff Moller}; el
autor real de \emph{Numerical Computing with MATLAB} es \textbf{Cleve B.
Moler}. Se corrige por ser un dato verificable (fundador de MATLAB, autor
del libro con el ejemplo citado en clase), no una invención.
\end{notebox}

Lógica del backslash de juguete:

\begin{enumerate}
  \item Si $A$ es triangular inferior $\to$ sustitución hacia adelante directa.
  \item Si $A$ es triangular superior $\to$ sustitución hacia atrás directa.
  \item Si $A$ es simétrica $\to$ intentar Cholesky (más barato que LU).
  \item Si nada de lo anterior aplica $\to$ calcular la LU (eliminación con
  pivoteo parcial) y resolver como en la Sección 5.
\end{enumerate}

\begin{algorithm}[H]
\SetAlgoLined
\KwEntrada{Matriz $A \in \mathbb{R}^{n\times n}$}
\KwSalida{$A$ sobrescrita: parte triangular inferior estricta $=L-I$, parte
triangular superior (con diagonal) $=U$; vector de permutación $p$}
$p \gets (1,2,\dots,n)$\;
\For{$k \gets 1$ \KwTo $n-1$}{
  $j^\ast \gets \operatorname*{argmax}_{j \geq k} |A_{jk}|$\;
  \If{$j^\ast \neq k$}{
    intercambiar filas $k$ y $j^\ast$ de $A$\;
    intercambiar $p_k \leftrightarrow p_{j^\ast}$\;
  }
  \If{$A_{kk} \neq 0$}{
    \For{$j \gets k+1$ \KwTo $n$}{
      $A_{jk} \gets A_{jk} / A_{kk}$ \tcp*{multiplicador $L_{jk}$, guardado in-place}
      \For{$i \gets k+1$ \KwTo $n$}{
        $A_{ji} \gets A_{ji} - A_{jk}\cdot A_{ki}$\;
      }
    }
  }
}
\caption{Descomposición LU con pivoteo parcial (versión de juguete, almacenamiento in-place, según Moler)}
\end{algorithm}

\begin{notebox}[Detalle de implementación]
Los multiplicadores se guardan directamente en la parte de $A$ que la
eliminación va dejando en ceros lógicamente: no hace falta memoria extra
para $L$ y $U$ por separado. Al terminar, la parte triangular inferior
estricta de la matriz sobrescrita son los multiplicadores ($L-I$) y la
parte triangular superior (con diagonal) es $U$.
\end{notebox}

\section{Matrices dispersas (introducción)}

En muchas aplicaciones (redes/grafos, sistemas estructurales como el
ejemplo del puente) las matrices tienen muchos ceros: una variable no está
directamente relacionada con todas las demás, solo con sus vecinas. Estas
se llaman \textbf{matrices dispersas} (sparse).

\begin{itemize}
  \item Almacenar una matriz dispersa como si fuera densa (los $n^2$
  números, incluyendo todos los ceros) es un desperdicio de memoria.
  \item En Octave, el comando \texttt{sparse} permite guardar solo las
  coordenadas y valores de las entradas no nulas. Ejemplo de clase: una
  matriz $50\times 50$ densa ocupa 20\,000 bytes; su versión dispersa, unas
  8 veces menos.
  \item La eliminación gaussiana (con o sin pivoteo) siempre funciona
  mientras $A$ sea no singular —es \enquote{la red de seguridad}— pero es
  $O(n^3)$ sin aprovechar ninguna estructura. Para matrices dispersas
  existen métodos más eficientes, que el curso no desarrolla en esta clase.
\end{itemize}

\begin{notebox}
Continúa la clase siguiente retomando este tema y avanzando hacia el
análisis de cuán bueno es, en la práctica, el método de eliminación con
pivoteo parcial (normas de matrices y número de condición — Clases 7 y 8).
\end{notebox}

\end{document}
